\chapter{Electron--ion collisions and field ionization}
\label{ch:coll}

\section{Electron-impact ionization}

With the electron cloud on, its electrons ionize the atoms they reach (\opt{mcmd-collisions = 1}, the default).
Every configuration can lose one electron from each occupied shell, to any final configuration that has data.
The rate of a channel $i\to f$ at atom $k$ is
\begin{equation}
  w_{i\to f} = n_e(\vect x_k)\,\langle\sigma v\rangle_{i\to f}(T),
\end{equation}
with the cloud's temperature $T$ and its density at the atom, interpolated from the last solve,
$n_e(\vect x_k) = \varsigma\sum_{\vect o}w_{\vect o}\,n_{\vect i_0+\vect o}/\Delta V$ (cloud-in-cell weights).
An ion deep in the quasi-neutral core sees a high density and ionizes quickly; one in the depleted outer sheath sees
almost none.
The collisional channels join the atom's other transitions in the Monte Carlo step, with the rates frozen for the
step, and the energy is booked as in Section~\ref{sec:energy} (item~3).
The transition energy is always $\Delta E = T_i - T_f$ from the total energies.

\subsection{Rate coefficients: Lotz (built in) or from files}

\opt{mcmd-coll-data} chooses where $\langle\sigma v\rangle$ comes from.
\begin{description}
  \item[\code{lotz}] (default): Lotz's Maxwellian formula~\cite{lotz1968} (with $b=0$),
    \begin{equation}
      \langle\sigma v\rangle = 6.7\times10^{-7}\,a\,q\,T^{-3/2}\,\frac{E_1(P/T)}{P/T}\ \ \mathrm{cm^3\,s^{-1}},\qquad
      a = 4.5\times10^{-14}\ \mathrm{cm^2\,eV^2},
    \end{equation}
    $q$ the occupation of the shell, $P = \Delta E$ (eV), $E_1$ the exponential integral.
    No data file is needed.
  \item[\code{files}]: from \file{collisional\_ionization\_table\_X.txt} if present, a table of
    $\langle\sigma v\rangle$ on a temperature grid that any atomic code can supply; else from the Sampson fits in
    \file{collisional\_parameters\_X.txt} (CRETIN models):
    \begin{equation}
      \langle\sigma v\rangle = 1.0899\times10^{-6}\,\frac{s_0\,q(b)}{\Delta E\sqrt T}\,e^{-b},\quad
      q(b) = c_0 + c_1\sqrt{\pi b}\,\mathrm{erfcx}\sqrt b + (c_2b+c_3)\,e^bE_1(b),\quad b = \frac{\Delta E}{T}.
    \end{equation}
    Tables are interpolated linearly in $\ln T$ in the smooth quantity $\ln\langle\sigma v\rangle + \Delta E/T$ and
    held at the grid edges outside them.
\end{description}
For H, C, N, O and S Lotz gives the same transitions as the CRETIN fits, with coefficients 13--25\,\% higher
between 10 and 300\,eV.
On lysozyme at 600\,eV this changes the ion kinetic energy by $-2.5$\,\% and nothing else beyond the seed noise.
A \file{.tpr} older than format 79 runs as \code{files}.

\section{Three-body recombination}

The inverse of collisional ionization, two cloud electrons colliding at an ion with one of them captured, follows by
detailed balance with the statistical weights $g$:
\begin{equation}
  w_{f\to i} = n_e^2\,\langle\sigma v\rangle_{i\to f}\,\frac{g_i}{g_f}\,1.656\times10^{-22}\,T^{-3/2}\ \ \mathrm{s^{-1}}.
\end{equation}
It is off by default (\opt{mcmd-coll-recombination = 0}) and \textbf{not fully validated}: in the classical
continuum picture it recaptures far too many electrons in the dense, cool cloud late in the explosion.
\opt{mcmd-coll-h-recombination} additionally allows recombination onto hydrogen (off: a proton stays a proton).

\section{Collisional excitation}

With \opt{mcmd-coll-excitation = 1} the cloud also excites atoms between configurations with the same number of
electrons, and de-excites them by detailed balance ($\langle\sigma v\rangle_{\mathrm{down}} =
\langle\sigma v\rangle_{\mathrm{up}}\,g_{\mathrm{lower}}/g_{\mathrm{upper}}$).
The rates come from \file{collisional\_excitation\_table\_X.txt} or the CRETIN \code{colex2} fits in
\file{collisional\_excitation\_X.txt} (Chapter~\ref{ch:data}).
It is off by default and \textbf{not fully validated}: it produces too much early collisional ionization out of excited
configurations.

\section{Field ionization}

In a strongly charged sample the collective field can lower an atom's barrier below its outermost binding energy.
With \opt{mcmd-field-ionization = 1} an atom is stripped while the field at it exceeds the Bethe over-the-barrier
threshold,
\begin{equation}
  F \ge \frac{I_p^2}{4\kc Z},
\end{equation}
$I_p$ the binding energy of its outermost electron and $Z$ its charge after the step.
$F = |\nabla P_L|$ is the field of the collective potential $P_L = \kc G_L*(q_{\mathrm{ion}}-n_e)$, with every
charge spread over $L$ (\opt{mcmd-field-ionization-smoothing}, 0.5\,nm; at least the grid's $s$) and the atom's own
charge left out, so neighbours closer than $\sim L$ are smeared: their barrier lowering is the charge-transfer
channel's.
It is computed on the grid after every solve, by direct summation where there is no grid, during the pulse only
(from 6$\sigma_t$ before to 4$\sigma_t$ after the peak, stopping once no atom has been within a factor 2 of its
threshold for 5 updates).
Every field ionization is logged in \file{pbfield\_events.txt}.

Field ionization does not change the explosion: on ubiquitin at 600\,eV it strips $\sim$900--1000 atoms per run early in
the pulse, and about as many fewer are ionized by electron impact and charge transfer; the charge, the escape and the
ion directions are unchanged.
It is off by default.
